\documentclass[12pt,a4paper]{article}
\usepackage[utf8]{inputenc}
\usepackage{geometry}
\usepackage{fancyhdr}
\usepackage{hyperref}

% Geometría de la página
\geometry{
 a4paper,
 total={170mm,257mm},
 left=20mm,
 top=20mm,
}

% Encabezado y pie de página
\pagestyle{fancy}
\fancyhf{}
\fancyhead[C]{Propuesta de Webinar: Seguridad en Criptoactivos}
\fancyfoot[C]{\thepage}
\renewcommand{\headrulewidth}{0.4pt}
\renewcommand{\footrulewidth}{0.4pt}

% Título
\title{
    \vspace{2cm}
    \textbf{Propuesta de Webinar Gratuito} \\
    \large \textit{Bitcoin: Orígenes, Funcionamiento y Principios de Seguridad}
    \vspace{1.5cm}
}
\author{Ing. Daniel Velarde Kubber}
\date{La Paz, 13 de junio de 2025}

\begin{document}

\maketitle
\thispagestyle{empty}
\clearpage

\section*{1. Detalle General de la Propuesta}

\begin{itemize}
    \item \textbf{Título del Webinar:} Bitcoin: Orígenes, Funcionamiento y Principios de Seguridad.
    \item \textbf{Propuesta dirigida a:} Universidad Real de la Cámara Nacional de Comercio.
    \item \textbf{Ponente:} Ing. Daniel Velarde Kubber - Ingeniero Industrial y de Sistemas, entusiasta en Data y Blockchain.
    \item \textbf{Fecha propuesta:} 13 de junio de 2025.
    \item \textbf{Ciudad:} La Paz, Bolivia.
    \item \textbf{Modalidad:} Virtual (a través de la plataforma que disponga la universidad).
    \item \textbf{Duración Total:} 1 hora (45 minutos de exposición y 15 minutos para preguntas y respuestas).
    \item \textbf{Costo:} Gratuito.
\end{itemize}

\section*{2. Objetivo del Webinar}

Ofrecer a la comunidad universitaria una introducción sólida y desmitificada a Bitcoin. El webinar tiene como fin explicar qué es, cómo funciona su red descentralizada y cuáles fueron sus orígenes. Se pondrá un énfasis crucial en los aspectos de seguridad, los riesgos asociados y las prácticas necesarias para interactuar con el primer criptoactivo de forma segura y responsable.

\section*{3. Temario y Contenido (Duración: 45 minutos)}

\begin{enumerate}
    \item \textbf{El Origen de Bitcoin (10 min)}
    \begin{itemize}
        \item El contexto pre-Bitcoin: La crisis financiera de 2008 y la necesidad de una alternativa.
        \item Satoshi Nakamoto y el Whitepaper de Bitcoin: "Un Sistema de Efectivo Electrónico Peer-to-Peer".
        \item El nacimiento del primer bloque (Bloque Génesis).
    \end{itemize}
    
    \item \textbf{¿Qué es Bitcoin y Cómo Funciona? (15 min)}
    \begin{itemize}
        \item Bitcoin como dinero digital descentralizado.
        \item Conceptos clave: Blockchain, minería (Prueba de Trabajo - Proof of Work), nodos y transacciones.
        \item La política monetaria de Bitcoin: Emisión finita (21 millones) y "halving".
        \item ¿Por qué se le considera "oro digital"?
    \end{itemize}
    
    \item \textbf{Seguridad en la Red Bitcoin (5 min)}
    \begin{itemize}
        \item La criptografía como pilar de la seguridad.
        \item Inmutabilidad y resistencia a la censura gracias a la descentralización.
    \end{itemize}
    
    \item \textbf{Riesgos y Uso Responsable (15 min)}
    \begin{itemize}
        \item \textbf{Riesgos principales:} Volatilidad de precios, ataques de ingeniería social, y errores de usuario.
        \item \textbf{La Responsabilidad de la Autocustodia:}
        \begin{itemize}
            \item "Not your keys, not your coins": La diferencia entre tener Bitcoin en un exchange y en una billetera propia.
            \item Tipos de billeteras: Hardware (frías), software (calientes) y de papel.
            \item Mejores prácticas para la gestión de claves privadas y frases semilla.
        \end{itemize}
        \item Cómo realizar transacciones de forma segura y verificar su estado en un explorador de bloques.
        \item Consejos para evitar estafas y proteger tus activos.
    \end{itemize}
\end{enumerate}

\section*{4. Sesión de Preguntas y Respuestas (15 minutos)}

Al finalizar la presentación, se abrirá un espacio para que los asistentes puedan realizar preguntas y resolver sus dudas directamente con el ponente.

\vspace{2cm}

\begin{center}
    Agradezco de antemano su tiempo y consideración. \\
    \vspace{0.5cm}
    \textbf{Ing. Daniel Velarde Kubber} \\
    \textit{Entusiasta en Data y Blockchain}
\end{center}

\end{document}